% !TEX root = ../thesis.tex

\chapter{Triton 算子调用的等价性验证问题定义}
\label{chap:operator-verification}

\section{研究对象与验证边界}
\label{sec:opv-scope}

本章定义一次 Triton 算子调用的形式化验证问题，为验证案例的表达与验证工具的实现提供依据。
研究对象是非原地、输入与输出不别名、仅产生一个逻辑输出张量的调用变体，
且候选实现通过一次完整的 kernel launch 完成计算。
该范围面向 ntops 等算子集合中满足上述条件的调用变体；
具体实现是否满足限制，需要逐个检查，不能仅根据算子名称判定。

验证单元包含一次 launch 中的全部 program instance，
以及构造该调用所需的输入映射、存储布局、编译期参数和 launch grid 生成规则。
单个 program、单条存储指令或一次符号执行路径均不是完整的验证单元。
一个案例可以量化多个形状，从而描述一族调用；每个输入实例对应其中一次完整调用。

第一阶段进一步要求各 program 只读取共同输入，独立完成计算，并写入互不重叠的输出区域。
program 内部可以包含向量运算与归约；跨 program 通信、全局原子累加式协作归约、
多次 launch 的组合、原地更新和额外辅助输出暂不属于本章的问题范围。
其中，program 之间写入不重叠以及整个输出被完整产生，均为待验证性质，不能作为未经检查的假设。

\section{算术模型与共同输入契约}
\label{sec:opv-contract}

\subsection{显式算术语义}

记 $\mathcal{M}$ 为案例采用的语义模型。
该模型明确规定数据运算、索引运算、类型转换及相关原语的含义。
第一阶段将浮点计算解释为数学实数上的运算，并要求模型内的精确相等。
涉及整数时，也必须明确所选整数语义；数学整数与具有固定位宽的整数不能隐式混用。
本章的参数化定义不预先替所有整数操作选择同一种解释。

实数抽象中的等价性不自动推出实际浮点执行的逐位相等或误差界。
例如，加法在实数模型中满足结合律，而实际浮点加法可能因舍入而不满足结合律。
因此，形式化结论应随所选模型一起报告；数值容差测试用于补充检查实际浮点行为。
对未被模型定义的操作，不能任意替换其语义后仍宣称证明了原调用。

\subsection{逻辑输入与契约}

\begin{definition}[逻辑输入与输入契约]
  逻辑输入 $x$ 包含形状、标量参数、输入张量的逻辑元素，
  以及案例显式声明的存储与布局信息。
  输入契约是这些量上的谓词 $P(x)$，对应的输入域为
  \begin{equation}
    \mathcal{D}_P = \{x \mid P(x)\}.
    \label{eq:opv-domain}
  \end{equation}
  reference 与 candidate 共用该逻辑输入域。
\end{definition}

张量元素在契约允许的范围内全称量化。
形状与标量参数可以由案例固定，也可以在显式约束下全称量化。
例如，固定矩阵乘法的 $M=64$、$N=128$、$K=32$，
仍然要求结论对所有满足契约的矩阵元素成立；具体形状不意味着只验证某个示例张量。
符号形状案例则同时量化允许的形状和每种形状下允许的张量取值。

\begin{definition}[案例有效性]
  输入契约必须可满足，即
  \begin{equation}
    \exists x\; P(x),
    \qquad\text{等价地，}\mathcal{D}_P\ne\varnothing.
    \label{eq:opv-satisfiable}
  \end{equation}
\end{definition}

契约可满足性独立于输出等价性检查。
若同时要求 $K=30$ 与 $K\bmod32=0$，输入域为空，案例应被判为无效。
不能利用蕴含式前件恒假而获得的真命题，报告有意义的验证成功。
若尚不能确定契约可满足，也不能报告完整验证成功。

\subsection{共同输入与双方调用映射}

记 $b_R$ 与 $b_C$ 分别为 reference 和 candidate 的调用映射。
两种映射使用同一个 $x$，构造各自需要的参数及存储表示。
双方执行使用独立状态，输入的逻辑值一致。
形状兼容性属于共同契约；指针、stride、编译期常量和 grid 等如何表示这些逻辑量，
属于调用映射。映射必须保持约定的逻辑输入和输出接口。

实现特有的前置要求不能作为某一侧的私有假设，隐式缩小共同输入域。
若 candidate 要求谓词 $Q_C(x)$，应证明
\begin{equation}
  \forall x,\quad P(x)\Longrightarrow Q_C(x),
  \label{eq:opv-requirement}
\end{equation}
或者显式建立契约包含 $Q_C$ 的受限案例。
后一种处理改变了所证明的输入范围，必须在案例中可见。

\section{参考函数与候选执行}
\label{sec:opv-semantics}

\subsection{参考输出函数}

记 $\mathcal{I}_Y(x)$ 为约定输出张量的逻辑索引集合，
$\mathcal{V}_{\mathcal{M}}$ 为所选模型中的元素值域。
reference 经调用映射及语义解释后给出参考输出函数 $R_{\mathcal{M}}$。
对每个 $x\in\mathcal{D}_P$，其输出具有类型
\begin{equation}
  R_{\mathcal{M}}(x):
  \mathcal{I}_Y(x)\longrightarrow\mathcal{V}_{\mathcal{M}}.
  \label{eq:opv-reference}
\end{equation}
该记号已包含 $b_R$，后文不再重复书写调用映射。

reference 必须在整个契约输入域上有定义，记为
\begin{equation}
  \forall x,\quad
  P(x)\Longrightarrow\operatorname{Def}_R(x).
  \label{eq:opv-reference-defined}
\end{equation}
例如，涉及除法时，应保证参考表达式在允许输入上不存在未定义的除零。
reference 不受单次 launch 限制；它可以通过数学表达式或具有明确语义的程序给出。
本章将其作为功能依据，证明候选实现符合该依据。
reference 是否准确表达开发者期望的算子功能，属于规格本身的正确性问题。

\subsection{候选调用与执行合法性}

候选程序、调用映射 $b_C$ 和语义模型共同确定一次完整调用。
记 $\mathcal{E}_C(x)$ 为模型允许的完整执行集合。
该集合必须覆盖模型允许的执行顺序，不能通过预先删除出错执行来定义。
无法正确构造调用、发生错误或无法完成的行为不能作为成功结果。

记 $\operatorname{Legal}_C(x)$ 为候选调用的执行合法性义务。
它要求调用可以正确构造，存在模型执行，且模型允许的全部执行均正常完成，
遵守所选模型的内存与并行执行规则，并满足本章的访问限制。
这些限制包括：输入保持只读；全局内存读取仅来自声明的输入；写入落在声明的输出存储范围内；
不存在越界、未定义操作或模型禁止的数据竞争。
输入、输出的物理表示必须实现约定的逻辑接口，输出存储不得与输入别名。
仅比较数值相等不能替代这些义务。

\subsection{输出所有权与完整性}

设 $\mathcal{G}(x)$ 为该调用的 program 标识集合。
对正常执行 $e\in\mathcal{E}_C(x)$，
记 $W_g(x,e)\subseteq\mathcal{I}_Y(x)$ 为 program $g$ 写入的逻辑输出元素集合。
通过输出的存储映射，将物理写入对应到逻辑索引；该映射必须保证不同逻辑输出元素具有独立存储。

各 program 的写入集合应构成输出索引集合的一个划分，允许其中的集合为空：
\begin{align}
  W_g(x,e)\cap W_h(x,e)&=\varnothing,
  &&g,h\in\mathcal{G}(x),\ g\ne h,
  \label{eq:opv-disjoint}\\
  \bigcup_{g\in\mathcal{G}(x)}W_g(x,e)&=\mathcal{I}_Y(x).
  \label{eq:opv-coverage}
\end{align}
式\eqref{eq:opv-disjoint}规定不同 program 的输出所有权不重叠；
式\eqref{eq:opv-coverage}规定没有应当产生却被遗漏的输出元素。
这些条件对所有契约输入和所有模型执行均需成立，统称为 $\operatorname{Partition}_C(x)$。
若存在非正常执行，则该义务不成立。
同一 program 内部的写入仍须满足执行合法性；集合记号本身不证明内部不存在竞争。

当上述义务成立时，记 $Y_C(x,e)$ 为完整执行结束后的逻辑输出张量。
可观察结果是 $\mathcal{I}_Y(x)$ 上所有元素的最终值，
不包含未被接口声明为结果的局部中间值或存储填充区域。
对填充区域不作结果比较，并不允许越过声明的合法访问范围。

\section{等价性验证任务}
\label{sec:opv-problem}

\begin{definition}[验证案例]
  一个验证案例由共同输入契约 $P$、reference、candidate、
  双方调用映射、显式语义模型 $\mathcal{M}$，
  以及固定参数与量化域的声明组成。
  这些内容共同确定一个待证命题。
  同一算子可以具有多个案例；改变形状实例化、输入限制或语义模型，可能改变待证命题。
\end{definition}

\begin{definition}[契约域上的调用等价性]
  一个案例通过完整验证，当且仅当以下条件全部成立：
  \begin{enumerate}
    \item 案例有效：式\eqref{eq:opv-satisfiable}成立。
    \item reference 在整个契约域上有定义：式\eqref{eq:opv-reference-defined}成立。
    \item candidate 在整个契约域上满足执行合法性和输出划分义务：
      \begin{equation}
        \forall x,\quad P(x)\Longrightarrow
        \bigl(\operatorname{Legal}_C(x)\land
        \operatorname{Partition}_C(x)\bigr).
        \label{eq:opv-candidate-obligations}
      \end{equation}
    \item 在上述定义性和合法性得到保证后，所有完整执行的输出与参考输出逐元素相等：
      \begin{equation}
        \begin{aligned}
          \forall x\in\mathcal{D}_P,\quad
          \forall e\in\mathcal{E}_C(x),\quad
          \forall i\in\mathcal{I}_Y(x),\qquad\\
          Y_C(x,e)[i]=R_{\mathcal{M}}(x)[i].
        \end{aligned}
        \label{eq:opv-equivalence}
      \end{equation}
  \end{enumerate}
\end{definition}

上述定义将契约有效性、执行合法性与结果等价性分别表述。
存在非法执行或缺失输出时，不能只比较其余已产生的结果并报告成功。
本章讨论的是共同契约域上的调用等价性，不对契约外的行为作出保证，
也不将其直接推广为一般部分程序的等价性或精化关系。

\begin{proposition}[可观察结果的确定性]
  若案例通过完整验证，则对任意 $x\in\mathcal{D}_P$，
  candidate 的所有模型执行均产生相同的逻辑输出。
\end{proposition}
\begin{proof}
  由式\eqref{eq:opv-equivalence}，任意两个执行 $e_1,e_2\in\mathcal{E}_C(x)$
  在每个逻辑输出索引 $i$ 上均等于 $R_{\mathcal{M}}(x)[i]$，
  故二者的输出逐元素相等。
\end{proof}

\section{矩阵乘法示例}
\label{sec:opv-matmul}

在实数模型中，令
\begin{equation}
  \symbf{A}\in\mathbb{R}^{M\times K},\qquad
  \symbf{B}\in\mathbb{R}^{K\times N},\qquad M,N,K>0.
\end{equation}
共同输出索引集合为
\begin{equation}
  \mathcal{I}_Y(x)=
  \{(i,j)\mid 0\le i<M,\ 0\le j<N\},
\end{equation}
参考函数为
\begin{equation}
  R_{\mathcal{M}}(x)[i,j]
  =\sum_{k=0}^{K-1}A_{ik}B_{kj}.
  \label{eq:opv-matmul-reference}
\end{equation}

候选调用采用正整数块大小 $B_M,B_N$，其 grid 为
\begin{equation}
  \mathcal{G}(x)=
  \{0,\dots,\lceil M/B_M\rceil-1\}
  \times\{0,\dots,\lceil N/B_N\rceil-1\}.
\end{equation}
program $(p,q)$ 负责的输出集合为
\begin{equation}
  \begin{aligned}
    W_{p,q}(x)=\{(i,j)\in\mathcal{I}_Y(x)\mid{}&
      pB_M\le i<(p+1)B_M,\\
      &qB_N\le j<(q+1)B_N\}.
  \end{aligned}
\end{equation}
这里的集合描述是实现应满足的目标；验证仍须证明实际写入与其一致。
每个 program 独立完成所负责元素沿 $K$ 维的全部归约。
由区间划分可检查输出所有权不重叠与覆盖性，
还需独立检查输入读取和尾块屏蔽是否合法，最后证明输出满足式\eqref{eq:opv-matmul-reference}。

若实现要求 $K\bmod32=0$，则应从共同契约推导该条件，
或者将其明确加入受限案例。
固定形状与符号形状案例可采用相同的参考函数，但对应的全称量化范围不同。
这使验证案例同时表达算子的功能、适用输入域和具体调用方式，而不依赖测试样本来隐式定义命题。

\section{结论的适用范围}

本章定义给出验证工具应当处理的完整数学任务，并不意味着某一实现已经具备全部证明能力。
求解超时、操作未建模或某项义务尚未完成，均不足以得出完整验证成功。
实际报告应明确采用的语义模型、契约输入域与已完成的证明义务。
在该范围内，形式化验证建立模型中的全称结论，
数值测试则补充检查实际浮点实现的行为；二者共同服务于高性能算子开发中的正确性检查。
